\section{Introduction}
Photovoltaic maximum power point tracking (MPPT) controls the electrical operating point of a source whose current--voltage characteristic depends on irradiance, temperature and module configuration. Its objective is to extract available energy over time, including environmental transitions and load disturbances. This objective differs from predicting an instantaneous electrical quantity accurately. An estimator can reduce its current residuals without locating the maximum of the power curve more accurately. Establishing whether online learning improves MPPT therefore requires both a predictive assessment and a closed-loop energy comparison.

The review by Esram and Chapman \cite{esram} organizes MPPT techniques according to their operating principles, measurement requirements and implementation characteristics. Subudhi and Pradhan \cite{subudhi} extend this comparison to control strategies and practical implementation. More recent classifications retain conventional, intelligent, optimization-based and hybrid families \cite{review2023}. These categories describe different mechanisms rather than a universal performance hierarchy. Their comparison depends on the source, sensing chain, converter, sampling interval and evaluation profile. In this study, conventional feedback provides the starting point for examining what a learned model adds to the control loop.

\subsection{Perturb and observe and incremental conductance}
Perturb and observe (P\&O) changes an operating variable and uses the resulting power difference to select the next perturbation direction. With a fixed step, a larger displacement can shorten approach to the maximum but increases the excursion around it. A smaller displacement reduces the steady-state excursion at the cost of slower movement across the characteristic. Femia et al. \cite{femia} analyze the choice of perturbation amplitude and observation interval in relation to plant behavior. Their treatment makes these design parameters part of the controller definition, not incidental implementation details. An observed power change may also contain a contribution from changing irradiance, so it cannot always be attributed to the preceding control action.

The comparative experiments of Hohm and Ropp \cite{hohm} show why classical baselines should not be treated as intrinsically weak. They compare optimized algorithms on common MPPT hardware and find that properly configured P\&O remains competitive. This matters when assessing a neural controller with prior information about the source. An apparent improvement over an arbitrary fixed step does not establish superiority over the classical method as a family. Common instrumentation and documented controller parameters are necessary for an interpretable comparison.

Incremental conductance (INC) uses the stationary-power condition, expressed through the relationship between incremental and instantaneous conductance. Hussein et al. \cite{hussein} provide an early formulation for tracking under changing atmospheric conditions. In a sampled implementation, finite differences approximate the characteristic slope. Voltage resolution, current noise and tolerances around the zero-slope condition consequently influence decisions. INC replaces the explicit power-perturbation decision with a slope-related test, but it does not remove the measurement or timing constraints of the feedback loop.

Sera et al. \cite{sera} directly examine P\&O and INC from mathematical and experimental perspectives. Their analysis demonstrates equivalence under the formulations and conditions studied, challenging the assumption that INC must outperform P\&O. The implication for the present comparison is methodological: observed differences must be interpreted through the implemented rules, step sizes, dead zones and data acquisition. Algorithm names alone do not explain a tracking-efficiency difference.

\subsection{Adaptive steps and the transition to hybrid control}
Variable-step methods address the compromise between convergence and steady-state oscillation by changing perturbation magnitude with operating conditions. Liu et al. \cite{liu} develop a variable-step INC method. Piegari and Rizzo \cite{adaptive} investigate adaptive P\&O, and Piegari et al. \cite{adaptive2015} examine perturbation selection and controller design with experimental testing. These works provide the basis for treating step adaptation as a design variable in its own right. They do not imply that one proportionality coefficient or slope rule transfers unchanged between plants.

Adaptation by a prescribed rule should also be distinguished from learning a model. A slope-dependent step can vary at every sample while its rule and parameters remain fixed. A neural predictor can use fixed weights and still generate a changing reference. Online learning requires observations to update the parameters of that predictor or policy. These distinctions become necessary when interpreting a controller that combines a pretrained network, a tracking regulator and a conventional fallback.

Classical and intelligent MPPT are not mutually exclusive categories. Zainuri et al. \cite{zainuri} combine adaptive P\&O with fuzzy control. Dawahdeh et al. \cite{hybrid2024} augment P\&O with a neural network to estimate a duty-cycle change during abrupt irradiance variations. Such architectures assign separate roles to a learned or rule-based component and the underlying feedback search. They motivate a component-level evaluation: a hybrid controller can recover after a disturbance because of its conventional branch, its learned reference, or both. Recovery by the complete system does not identify which component produced the benefit.

\subsection{Intelligent MPPT: representations, learning and implementation}
The use of neural networks for photovoltaic tracking predates current edge-computing platforms. Hiyama et al. \cite{hiyama} evaluate a neural-network-based real-time maximum-power tracking controller in 1995. Mellit and Kalogirou \cite{mellit} subsequently review artificial-intelligence applications across photovoltaic systems. Salam et al. \cite{salam} focus on soft-computing MPPT and discuss tracking, complexity and hardware implementation. This history places the contribution of a new study in its specific representation, learning mechanism and evidence, rather than in the use of artificial intelligence alone.

Elobaid et al. \cite{elobaid} classify neural MPPT according to input variables and controller structure. Electrical measurements, environmental measurements and mixed inputs require different sensing arrangements. A network can estimate a reference voltage directly, supply a control correction, or represent a characteristic from which the operating point is calculated. The architecture therefore determines what training labels are required. Controller commands are appropriate labels for policy imitation only when their provenance and intended behavior are defined. Measurements of current at different voltages instead support characteristic identification. Neither dataset automatically supplies independently verified optimal commands.

Model-based neural MPPT is particularly close to the present approach. Zečević and Rolevski \cite{nnmodel} derive iterative tracking rules from a neural representation of photovoltaic output current. Their formulation separates characteristic approximation from the search for the maximum. Rizzo and Scelba \cite{rizzo} address rapidly variable shading with an ANN-based MPPT method. These examples show that the represented characteristic and its operating domain matter as much as the network architecture. A source with a single power maximum does not test the same search problem as an array with multiple maxima under partial shading.

Reinforcement learning provides another formulation in which an agent selects actions according to an objective defined through rewards. Giraldo et al. \cite{giraldo} evaluate a deep-Q-network in simulations and real photovoltaic experiments, including partial shading. This differs from updating a supervised current predictor from measured current. In the present implementation, the current residual drives learning and energy tracking is assessed afterward. Improvements in the learning target therefore need not produce improvements in the control objective. Introducing a reward-based policy would require a different training and validation design.

Recent implementation studies also distinguish forms of adaptation and levels of validation. Abouzeid et al. \cite{abouzeid} combine numerical and physical experimental evaluation of an ANN MPPT against conventional controllers. Amine et al. \cite{amine} use meteorological acquisition through a Raspberry Pi to support dataset expansion and incremental estimator updates. Their workflow differs from local output-layer updates triggered by each eligible electrical observation. Jamshed et al. \cite{jamshed} deploy a lightweight neural controller on an STM32 for controller hardware-in-the-loop assessment. Physical experimentation, physical-controller HIL and simulated firmware execution provide different evidence. Their reported efficiencies and execution times should not be ranked without accounting for those differences.

\subsection{Relationship to the digital-panel framework and study objective}
The digital source and sensing workflow originate from the framework of Suárez Suárez et al. \cite{suarez}. That study provides the basis for evaluating MPPT through firmware-level measurements with sensing limitations. The present work extends the controller investigation by introducing a neural characteristic model with bounded online output-layer updates. Source provenance is retained, and the added automatic load interface is distinguished from the original panel equations.

The project contains records labeled PO, POINC and POINCADAP, corresponding to P\&O, INC and adaptive P\&O in the available firmware. These records supply conventional reference behavior and expose the roles of perturbation direction, slope information and step magnitude. They are not three standardized names for separate neural models. In the current intelligent controller, bounded reference tracking follows these control considerations, while INC supplies fallback behavior outside the covered neural domain. The implemented network weights are fitted to calibrated current--voltage sweep measurements. Historical control commands are not asserted to be optimal training labels or to have been combined into the deployed weights.

The resulting research question is specific: does updating the current predictor during operation improve energy tracking relative to the same predictor with frozen parameters? The study addresses it through a paired online-versus-frozen ablation, conventional P\&O/INC/adaptive-P\&O references, and a common acquisition and actuation path. Predictive residuals, independent maximum-power references and energy integrals are evaluated separately. The documentary comparison with ten commercial modules characterizes the digital source's consistency and limitations. It does not replace an independent MPPT reference for that source.

Table~\ref{tab:literature} summarizes the literature-to-design connection. The intended contribution is an auditable separation of model identification and tracking behavior in an ESP32-oriented simulation workflow. The study does not claim a new general family of neural MPPT, replacement of physical-converter experiments, or an energetic advantage before the paired tests establish it.

\begin{table}[htbp]
\caption{Connection between prior MPPT research and the present evaluation. The rows synthesize mechanisms and design consequences, not a ranking of efficiencies across unrelated test benches.}
\label{tab:literature}
\centering\small
\begin{tabularx}{\textwidth}{p{0.23\textwidth}LL}\toprule
Research line and references & Established consideration & Consequence for this study\\\midrule
Comparative MPPT foundations \cite{esram,subudhi,hohm} & Method selection depends on implementation and test conditions. & Common sensing, actuation, reset and sampling for all variants.\\
P\&O and INC \cite{femia,hussein,sera} & Perturbation, slope estimation and timing affect observed behavior. & Document steps, tolerances and measurement resolution.\\
Variable-step MPPT \cite{liu,adaptive,adaptive2015} & Step adaptation addresses convergence and steady-state behavior. & Separate a changing control step from parameter learning.\\
Hybrid intelligent feedback \cite{zainuri,hybrid2024} & Fuzzy or neural components can augment classical search. & Attribute recovery to the active branch and compare with learning frozen.\\
Neural representations \cite{hiyama,elobaid,nnmodel} & Reference prediction, policy correction and characteristic identification are distinct tasks. & State the training target and retain its provenance.\\
Variable shading and learned policies \cite{rizzo,giraldo} & Global tracking and action learning require their own problem definitions. & Do not infer global-MPP capability from a single-maximum plant.\\
Recent implementation \cite{abouzeid,amine,jamshed} & Experimental, incremental-data and HIL workflows provide different evidence. & Report validation level and computational context explicitly.\\
Digital sensing framework \cite{suarez} & Instrumentation is part of the observable control loop. & Preserve the plant and audit scaling before attributing effects to learning.\\\bottomrule
\end{tabularx}
\end{table}


\subsection{Classical controllers and their role in the intelligent design}
The file labels PO, POINC and POINCADAP denote P\&O, INC and adaptive P\&O, respectively. In particular, the POINCADAP label is not interpreted as evidence of a combined P\&O--INC neural teacher. The present comparison uses the algorithm branches in the supplied firmware. Their conceptual foundations follow the perturbation, conductance and adaptive-step literature \cite{femia,hussein,liu,adaptive}.

P\&O retains a direction variable and reverses it when the observed power increment is negative. Its fixed perturbation is 1 percentage point of load. INC compares finite differences with the stationary-power condition
\begin{equation}
 \frac{dP}{dV}=I+V\frac{dI}{dV}=0,
 \qquad \frac{dI}{dV}=-\frac{I}{V}.
\end{equation}
The implemented branch uses voltage and current tolerances of 0.05 V and 0.02 A. When voltage changes exceed the voltage tolerance, the conductance comparison uses a threshold of 0.01 A/V. Load changes are bounded by the available 0--100\% command interval. Since increasing load reduces voltage in this plant, the sign of a voltage-directed correction must be mapped accordingly.

The adaptive P\&O branch retains the power-based direction rule and varies the load step according to
\begin{equation}
 a_n=\operatorname{clip}\left(0.08\left|\frac{\Delta P_n}{\Delta V_n}\right|,0.2,3.0\right),
 \qquad |\Delta V_n|>0.01\ \mathrm{V}.
\end{equation}
The minimum step is used otherwise. Steps are expressed in percentage points and the coefficient has the corresponding units per A. This is the particular rule evaluated here, not a claim to reproduce every adaptive method in the cited literature. Its fixed coefficients make it rule-based adaptation.

These classical mechanisms inform the intelligent design through signed reference correction, bounded increments, a near-reference dead zone and INC fallback. The neural component supplies a current model and a predicted reference rather than copying all classical decisions. Historical controller logs document trajectories and command behavior. The calibrated sweeps described next supply the current targets used to fit the deployed weights. This separation prevents controller commands from being mislabeled as measured maximum-power ground truth.

